\section{Background and Related Work}\label{sec:related}
This section introduces the datasets and evaluation metrics used in this study, followed by a
review of related work on MOT.

\subsection{Dataset Description}\label{sec:data}
Three publicly available video datasets, MOT17 \cite{milan2016mot16}, KITTI tracking
\cite{geiger2012kitti}, and UAVDT \cite{du2018uavdt}, are used to develop and evaluate the
proposed \scl{}. They differ in their objects, their cameras, and their recording conditions, and
they are used with different detectors. Using these datasets allows us to examine \scl{} with
well-matched detections, with other detectors, and on aerial videos that were not used during the
design. Table~\ref{tab:data} lists the evaluation cells; a cell is one combination of data,
detector, and tracker.

\subsubsection{MOT17}
MOT17 is a widely used benchmark for tracking people \cite{milan2016mot16}. It contains 14 video
sequences, 7 for training and 7 for testing, filmed in public places with static and moving
cameras. As is common, each of the seven training videos is split in time, and the second half,
called val-half, is used for evaluation \cite{zhang2022bytetrack}; it has 2,652 frames. On MOT17
the trackers use the published detections of the YOLOX-X detector \cite{ge2021yolox}, so every
tracker sees exactly the same boxes. A detector drops the boxes whose score is below a small value
before it passes them on; we call this value the emission floor. The two published MOT17 detection
sets have floors of 0.01 and 0.1. These detections fit the published trackers well, so MOT17 shows
whether \scl{} leaves a well-matched tracker unchanged.

To study the score problem directly, the published MOT17 scores were also changed before the
tracker by four monotone transforms, which change the score values but keep their order: cubing
($s^3$), halving ($0.5\,s$), and logit temperatures $T=2$ and $T=0.5$, $\sigma(z/T)$, where
$\sigma$ turns a logit back into a score. A temperature divides the logit $z$ of each score by $T$
before turning it back into a score: $T=2$ pulls all scores toward 0.5, and $T=0.5$ pushes them
toward 0 and 1. These transforms imitate the score changes that a new or retrained detector can
cause. The trackers kept their published thresholds.

\subsubsection{KITTI tracking}
KITTI tracking is a benchmark recorded from a car that drives in and around Karlsruhe, Germany
\cite{geiger2012kitti}. Its training set has 21 sequences with 8,008 frames, and its official
evaluation scores cars and pedestrians. In this study the trackers run on KITTI with two detectors
that differ from the MOT17 detector: YOLOv8n \cite{jocher2023yolov8}, a small and fast model, and
RT-DETR-L \cite{zhao2024rtdetr}, a larger transformer model. The trackers keep the thresholds that
were set for other detectors, so KITTI shows what happens when the detector is replaced.

\subsubsection{UAVDT}
UAVDT is a benchmark designed for object detection and MOT in videos from unmanned aerial
vehicles (UAVs), with a particular focus on vehicles \cite{du2018uavdt}. It consists of 100 video
sequences selected from more than 10 hours of video recorded at urban locations such as squares,
arterial streets, toll stations, highways, crossings, and T-junctions, at 30 frames per second and
a resolution of $1080\times540$ pixels. The sequences are recorded under different weather
conditions, flying altitudes, and camera viewing angles. The test set contains 20 sequences for the
MOT task, with 16,592 frames, and these 20 sequences are used in this study. They were not used
during the design, and ByteTrack runs on them with YOLOv8n and RT-DETR-L, restricted to their
vehicle classes.

The datasets serve different purposes in the evaluation. Development cells were looked at while
\scl{} was designed: MOT17 with four trackers, KITTI with two other detectors, and the score
transforms. The other cells test \scl{} on trackers and data that it had never seen. The two
predeclared external trackers, PD-SORT and Hybrid-SORT, were chosen and named before they were
run, and they were run once with the frozen configuration on the detections released by their
authors on MOT17 val-half. The held-out UAVDT test split was also run once with the frozen
configuration, after a written plan. To check that each external tracker was run correctly, our run
of the tracker alone is compared with the number reported by its authors and marked EXACT (within
0.2 HOTA points) or CLOSE (within 1.0 HOTA point, with a known cause that was not tuned).

\begin{table*}[t]
\caption{Evaluation cells and their role. A cell is one combination of data,
detector, and tracker. Development cells were looked at while the layer was
designed; predeclared and held-out cells were chosen in advance and run once
with the frozen layer.}
\label{tab:data}
\centering\footnotesize
\setlength{\tabcolsep}{4pt}
\begin{tabular}{@{}>{\raggedright\arraybackslash}p{29mm}>{\raggedright\arraybackslash}p{25mm}>{\raggedright\arraybackslash}p{42mm}>{\raggedright\arraybackslash}p{25mm}@{}}
\toprule
Data & Detector & Trackers & Role\\
\midrule
MOT17 val-half (7 sequences) & published YOLOX-X detections & ByteTrack (2 settings), OC-SORT, BoostTrack, SparseTrack & development\\
KITTI tracking training (21 sequences) & RT-DETR-L; YOLOv8n & ByteTrack and BoT-SORT (RT-DETR-L); OC-SORT (YOLOv8n) & development (detector transfer)\\
MOT17 val-half, score transforms & published YOLOX-X, transformed scores & ByteTrack (2 settings), OC-SORT, BoostTrack & development (robustness)\\
MOT17 val-half & released PD-SORT / Hybrid-SORT detections & PD-SORT, Hybrid-SORT & predeclared external\\
UAVDT test (20 sequences) & RT-DETR-L; YOLOv8n & ByteTrack & predeclared held-out\\
KITTI tracking training, also with halved scores & RT-DETR-L; YOLOv8n & ByteTrack (RT-DETR-L); OC-SORT (YOLOv8n) & development (baselines, ACT, and components)\\
\bottomrule
\end{tabular}
\end{table*}


\subsection{MOT Evaluation Metrics}\label{sec:metrics}
To evaluate \scl{}, six commonly used MOT metrics are considered: Higher Order Tracking Accuracy
(HOTA), Multiple Object Tracking Accuracy (MOTA), Identity F1 Score (IDF1), identity switches
(IDS), false positives (FP), and false negatives (FN). These metrics measure different aspects of
tracking performance, including detection accuracy, object identity consistency, and the ability to
track objects across consecutive frames \cite{dendorfer2021motchallenge,luiten2021hota}. They are
defined below, starting from the errors counted in each frame and ending with the scores that
combine them.

\emph{Detection errors.} The first metrics count the errors made frame by frame. In each frame
$t$, output boxes are matched one to one to ground-truth boxes with an intersection over union
(IoU) of at least 0.5. Let $\mathrm{GT}_t$ be the number of ground-truth boxes, $\mathrm{TP}_t$ the
number of matched pairs, $\mathrm{FP}_t$ the number of output boxes without a match (false
positives), and $\mathrm{FN}_t$ the number of ground-truth boxes without a match (false negatives,
or missed objects). Over all frames,
\begin{align}
\mathrm{GT}&=\sum_t \mathrm{GT}_t,\qquad \mathrm{TP}=\sum_t \mathrm{TP}_t,\label{eq:gt}\\
\mathrm{FP}&=\sum_t \mathrm{FP}_t,\qquad \mathrm{FN}=\sum_t \mathrm{FN}_t,\label{eq:fpfn}
\end{align}
so that $\mathrm{TP}=\mathrm{GT}-\mathrm{FN}$.

\emph{Identity switches.} The counts above treat each frame separately, so they do not show
whether an object keeps the same track over time. Identity switches measure this. Let
$\phi_t(i)$ be the output track matched to ground-truth object $i$ in frame $t$, and
$\phi_{\mathrm{prev}}(i)$ the output track to which object $i$ was last matched before $t$. An
identity switch is counted whenever these two tracks differ:
\begin{equation}
\mathrm{IDS}=\sum_t\sum_{i}\mathbf{1}\!\left[\phi_t(i)\neq\phi_{\mathrm{prev}}(i)\right],
\label{eq:ids}
\end{equation}
where the inner sum runs over the objects that are matched in frame $t$ and were matched before.
A tracker that follows more objects has more chances to switch identities, so IDS is read together
with FN.

\emph{MOTA and IDF1.} These two scores summarize the errors above in a single number each. MOTA
combines the three error types \cite{bernardin2008clear}, and IDF1 measures how long each object
keeps the correct identity \cite{ristani2016idf1}:
\begin{align}
\mathrm{MOTA}&=1-\frac{\mathrm{FN}+\mathrm{FP}+\mathrm{IDS}}{\mathrm{GT}},\label{eq:mota}\\
\mathrm{IDF1}&=\frac{2\,\mathrm{IDTP}}{2\,\mathrm{IDTP}+\mathrm{IDFP}+\mathrm{IDFN}}.\label{eq:idf1}
\end{align}
For IDF1, each ground-truth track is first matched to at most one output track over the whole
video. IDTP, IDFP, and IDFN are then the true-positive, false-positive, and false-negative boxes
under this track-level matching. Because MOTA counts every false positive and every missed object,
it reacts strongly to how many boxes reach the tracker. IDF1 depends more on whether each object
keeps one identity.

\emph{HOTA.} MOTA depends mostly on detection errors, and IDF1 mostly on identity errors. HOTA
balances detection and association in one score \cite{luiten2021hota}, and it is the main metric
of this paper. At a localization threshold $\alpha$, let $\mathrm{TP}_\alpha$, $\mathrm{FN}_\alpha$,
and $\mathrm{FP}_\alpha$ be the matched pairs, missed ground-truth boxes, and unmatched output boxes.
For a matched pair $c$, $\mathrm{TPA}(c)$ is the set of matched pairs with the same ground-truth
identity and the same output identity as $c$; $\mathrm{FNA}(c)$ is the set of boxes of the same
ground-truth object that are matched to another output identity or not matched; and
$\mathrm{FPA}(c)$ is the set of boxes of the same output track that are matched to another object
or not matched. Then
\begin{align}
\mathrm{DetA}_\alpha&=\frac{|\mathrm{TP}_\alpha|}{|\mathrm{TP}_\alpha|+|\mathrm{FN}_\alpha|+|\mathrm{FP}_\alpha|},\label{eq:deta}\\
\mathcal{A}(c)&=\frac{|\mathrm{TPA}(c)|}{|\mathrm{TPA}(c)|+|\mathrm{FNA}(c)|+|\mathrm{FPA}(c)|},\label{eq:ac}\\
\mathrm{AssA}_\alpha&=\frac{1}{|\mathrm{TP}_\alpha|}\sum_{c\in\mathrm{TP}_\alpha}\mathcal{A}(c),\label{eq:assa}\\
\mathrm{HOTA}&=\frac{1}{19}\sum_{\alpha}\sqrt{\mathrm{DetA}_\alpha\,\mathrm{AssA}_\alpha},\label{eq:hota}
\end{align}
where $\alpha$ runs over $0.05, 0.10, \dots, 0.95$. DetA measures how well the objects are found,
and AssA how well their identities are kept.

\emph{Benchmark protocols.} The metrics are computed with the TrackEval library (commit 12c8791)
and the official ground truth of each benchmark, following MOTChallenge practice
\cite{dendorfer2021motchallenge}. On KITTI tracking, HOTA is the official value averaged over cars
and pedestrians. On UAVDT, all vehicles form one class, and predictions inside the regions that the
dataset marks as ignored are removed.

\emph{Statistical comparison.} All comparisons are paired: each tracker runs twice on the same
sequences and detections, once alone and once with \scl{}, and nothing else changes. To judge
whether the difference between the two runs is reliable, paired bootstrap resampling is applied at
the sequence level \cite{efron1993bootstrap}. Resampling whole sequences keeps the dependence within
each sequence and gives intervals for the differences. Each resample draws the evaluation sequences
with replacement, 10,000 times with the fixed random seed 42. The same draw is used for both runs,
the pooled metrics are recomputed, and their difference is recorded. The 95\% interval is the
2.5--97.5 percentile range, and a difference is called significant when its interval does not
include zero. An interval that includes zero does not show that there is no difference; it shows
only that the tested sequences cannot separate the difference from zero. As a simple check that needs no statistics, we also count the sequences where HOTA
rose, stayed the same ($|\Delta\mathrm{HOTA}|<0.01$), or fell.

\emph{Units and direction.} HOTA, MOTA, and IDF1 are given in percent, and differences between two
runs in percentage points (points); relative changes are marked as such. IDS, FP, and FN are counts;
runtime is given in milliseconds (ms) per frame. Higher values are better for HOTA, MOTA, and IDF1,
and lower values are better for IDS, FP, FN, and runtime.

\subsection{MOT Literature Survey}\label{sec:survey}
MOT has been widely studied, with research focusing on better detectors, better association, and
better tracking under changing conditions. Three lines of this work relate directly to \scl{}.
Tracking-by-detection methods provide the trackers that \scl{} serves. Adaptive threshold methods
change thresholds inside a tracker and are the closest alternatives to \scl{}. Detector calibration
and adaptive inference methods change the detector scores or how a model runs. The following
subsections review these three lines in this order.

\subsubsection{Tracking-by-Detection Methods in MOT}\label{sec:hosts}
Tracking-by-detection is one of the most widely used approaches in MOT \cite{adzemovic2026survey}.
Objects are detected in each frame and then associated across frames to keep their identities.
SORT predicts where each tracked box will be in the next frame with a Kalman filter, a standard
motion predictor, and matches these predictions to the new detections by their overlap, using the
Hungarian method, which finds the best one-to-one matching \cite{bewley2016sort}. ByteTrack added a
second matching step that uses low-score detections to continue existing tracks
\cite{zhang2022bytetrack}, and many later trackers kept this two-stage design. Others, such as
OC-SORT \cite{cao2023ocsort}, use only one score threshold. All these trackers read detector scores
through fixed thresholds that were tuned for one detector, usually YOLOX \cite{ge2021yolox}. A recent
survey groups the tracking-by-detection methods into joint detection and embedding, heuristic,
motion-based, affinity-learning, and offline methods, and compares them across benchmarks
\cite{adzemovic2026survey}.

Many recent trackers improve their association with parts that are trained or designed for a
domain: small re-identification networks, which compare the appearance of objects
\cite{li2024fdbtrack}, learnable graph matching \cite{bao2024graphmatching,zhang2025sinkhorn},
result-guided feature extraction for faster tracking \cite{ding2025remot}, motion models for
low-frame-rate surface-vehicle videos \cite{liang2024motiontrack}, and image restoration by blind
super-resolution before tracking in adverse weather \cite{sun2026bsr}. These methods change or
retrain the tracker. \scl{} keeps both the tracker and the detector unchanged and acts only on what
passes between them.

\emph{Trackers and detectors used in this study.} Seven published trackers serve as hosts, that
is, as the trackers that \scl{} serves: ByteTrack \cite{zhang2022bytetrack}, BoT-SORT
\cite{aharon2022botsort}, OC-SORT \cite{cao2023ocsort}, SparseTrack \cite{liu2025sparsetrack},
BoostTrack \cite{stanojevic2024boosttrack}, PD-SORT \cite{wang2025pdsort}, and Hybrid-SORT
\cite{yang2024hybridsort}. Each host uses the authors' code and published configuration. ByteTrack
is used in two published settings: the official setting, with the thresholds that its authors used
on MOT17, and the library setting, with the default thresholds of the widely used Ultralytics
library, which also provides the BoT-SORT host. The detectors are YOLOX-X \cite{ge2021yolox},
YOLOv8n \cite{jocher2023yolov8}, and RT-DETR-L \cite{zhao2024rtdetr}. All hosts and detectors are
used without modification.

\subsubsection{Adaptive Threshold Selection in MOT}
Several studies adapt the thresholds of a tracker to changing conditions. Adaptive ByteTrack sets
the score threshold of ByteTrack from the detections of each frame \cite{ma2023adaptivebytetrack}.
AdapTrack adapts the thresholds of its matching step \cite{shim2024adaptrack}. ACT sets the high and
low thresholds of each frame from the first quartile of the scores and the number of candidates
\cite{act2026adaptive}. Other recent trackers use the score or the spread of each detection inside
the tracker. ConfTrack weights the Kalman update, the Kalman noise, and the matching cost by the
score of each box \cite{jung2024conftrack}. UTrack measures the spread of the boxes that
non-maximum suppression (NMS), the detector step that removes overlapping boxes of the same object,
merges, and uses it as the Kalman noise, so it needs the boxes before this step
\cite{solano2025utrack}. PlugTrack blends the Kalman filter with learned motion predictors through a
small network trained with labels \cite{kim2026plugtrack}. These methods are built into the tracker,
use only the current frame or each box alone, or need extra trained parts or the inside of the
detector.

\scl{} follows a different approach and differs in four ways. It uses past frames. It decides for
the whole video whether to act at all, and otherwise keeps the tracker's own settings. It works
through a host contract, four published settings of each tracker, on several published trackers
without changing them. And it is tested as one frozen configuration under score changes and on
external systems.

\subsubsection{Detector Confidence Calibration and Adaptive Inference}
\begin{table*}[t]
\caption{Comparison of \scl{} with recent adaptive methods for tracking (BT: ByteTrack; ReID: re-identification network; NMS: non-maximum suppression).}
\label{tab:related}
\centering\footnotesize
\setlength{\tabcolsep}{3pt}
\begin{tabular}{@{}>{\raggedright\arraybackslash}p{18mm}c>{\raggedright\arraybackslash}p{21mm}>{\raggedright\arraybackslash}p{19mm}>{\raggedright\arraybackslash}p{15mm}>{\raggedright\arraybackslash}p{16mm}>{\raggedright\arraybackslash}p{21mm}@{}}
\toprule
Method & Year & Adapted & Statistics & Extra trained part & Detector input & Where it works\\
\midrule
Adapt.\ BT \cite{ma2023adaptivebytetrack} & 2023 & high/low score split & current frame & none & boxes and scores & inside BT\\
ConfTrack \cite{jung2024conftrack} & 2024 & Kalman update, noise, and cost & score of each box & ReID & boxes and scores & new tracker\\
AdapTrack \cite{shim2024adaptrack} & 2024 & matching thresholds & current frame & ReID & boxes and scores & new tracker\\
UTrack \cite{solano2025utrack} & 2025 & Kalman noise & spread of boxes merged by NMS & none & boxes before NMS & added to existing trackers\\
ACT \cite{act2026adaptive} & 2026 & high and low thresholds & current frame & none & boxes and scores & inside the tracker\\
PlugTrack \cite{kim2026plugtrack} & 2026 & mix of Kalman and learned motion & track boxes & network trained with labels & boxes & added to two learned predictors\\
K{\"u}ppers \cite{kuppers2020detcalib} & 2020 & detector scores & labelled data & calibration map & boxes and scores & detector output\\
Chameleon \cite{jiang2018chameleon} & 2018 & pipeline configuration & profiling & none & whole pipeline & analytics pipeline\\
AC-MOT (ours, submitted) & 2026 & detector operating point & image and track cues & none & sets its settings & before the detector\\
\scl{} (this paper) & -- & candidates, scores, thresholds & past frames, no labels & none & boxes and scores & between detector and seven unchanged trackers\\
\bottomrule
\end{tabular}
\end{table*}

Scores do not mean the same thing for every detector. Detectors of different families, such as
YOLOX \cite{ge2021yolox}, YOLOv8 \cite{jocher2023yolov8}, YOLOv9 \cite{wang2024yolov9}, DETR
\cite{carion2020detr}, and RT-DETR \cite{zhao2024rtdetr}, give different score distributions. Their
confidence is not calibrated across models and conditions, and calibration methods for object
detection learn a mapping from labelled data \cite{kuppers2020detcalib}. A tracker tuned for one
detector therefore cannot assume that another detector's scores mean the same thing, and labels to
recalibrate them are rarely available after deployment. Image processing offers a tool that needs
no labels: Otsu's method picks a threshold from a histogram so that the two resulting groups are as
separate as possible \cite{otsu1979}. \scl{} applies it to the logits (log-odds) of the detector
scores of recent frames.

A wider group of methods adapts inference itself. Dynamic neural networks adapt their computation
to each input \cite{han2022dynamic}, and Chameleon measures again, from time to time, which settings
of a video analysis system work best \cite{jiang2018chameleon}. These methods learn or measure their
settings for each model. Our earlier work, AC-MOT (a manuscript submitted for publication), is a
controller that needs no training and sets the detector operating point (score threshold, NMS
threshold, and input resolution) while running, from image and track cues. Such a controller is
tied to the detector it was calibrated with, because a change in the detector's score scale does
not change the image. \scl{} reads the scores instead and acts on the tracker's input and
thresholds.

In summary, existing work either adapts thresholds or uses the score of each box inside one
tracker, calibrates scores with labelled data, or learns its settings for one model. A layer that
needs no training, adjusts itself from the detector's own scores, works with many unchanged
trackers, and keeps a well-matched tracker at its own accuracy is missing. \scl{} is designed to
fill this gap. Table~\ref{tab:related} summarizes the main differences between \scl{} and the
closest adaptive methods.

